\section{Artillery}

\begin{rulebox}{General Rule}
Each player controls a number of artillery pieces at the start of the game. These pieces are not units. They cannot move or participate in combat. However, each piece can be used to bombard (fire on) enemy units during a friendly artillery phase. Leaders and artillery pieces are never affected by a bombardment. Artillery pieces are never disrupted or eliminated. However, they can be captured and may be used to bombard enemy units.
\end{rulebox}

\ruleslabel{Procedure}

During his artillery phase, the active player can, at his option, bombard any one enemy unit with each artillery piece under his control. Artillery pieces under a player's control include those that started the game under that player's control and that are not currently captured by the opposing player, as well as all the opposing player's artillery pieces that are currently captured. For each bombardment, perform the following steps in order:

\begin{steps}
  \item Trace a path of hexes consisting of the fewest possible hexes in a straight line between the bombarding artillery piece (exclusive) and the enemy unit to be bombarded (inclusive). The result is the range. If any part of this path is blocked (7.3), the bombardment cannot take place, and the active player can choose another target unit (if one exists).
  \item Read across the row of numbers at the top of the Artillery Table until reaching the number equal to the range found in step 1. This is the range column for the bombardment.
  \item Roll a die, adding 1 to the roll if the target unit is in a close, train, or woods hex. Read down the possible die-roll results in the left-hand column of the Artillery Table until reaching the number that matches the modified die roll.
  \item Read across the row corresponding to the modified die roll until reaching the result in the range column found in step 2. This is the result of the bombardment.
  \item If the result is a Dd and the target unit is ordered, it is immediately disrupted; flip the unit to its back (disrupted) side. If the result is a dash, or if the result is Dd and the target unit is already disrupted, the bombardment has no effect. Regardless of the effect on the target unit, artillery pieces and leaders (of either side) are never affected by a bombardment.
\end{steps}

\ruleslabel{Cases}

\caseheading{7.1} The active player announces bombardments individually in any order he chooses.

Each bombardment is resolved immediately after being announced and before any other bombardment is resolved. Thus, it is possible to bombard an enemy unit any number of times during the same artillery phase, trying to disrupt a key unit with another bombardment should the previous bombardment fail to have any effect.

\caseheading{7.2} An enemy unit can be bombarded any number of times during an artillery phase. However, each artillery piece can bombard only one enemy unit.

\caseheading{7.3} Artillery pieces can only bombard enemy units that they can ``see.''

An artillery piece can see into and out of any type of terrain. However, artillery pieces cannot see through woods, close, town, train, hilltop, or plump hexes, or hexes containing artillery pieces or units of either side. These types of hexes are called blocking hexes because they block the line of sight to the target. \textsc{Exceptions:} See cases 7.4 and 7.5.

\caseheading{7.4} Artillery pieces in hilltop hexes that are firing into non-hilltop hexes ignore blocking hexes that are more than half the total distance in hexes from the hex occupied by the target unit.

\caseheading{7.5} Artillery pieces in hilltop hexes that are firing into other hilltop hexes ignore blocking hexes except for the plump hexes (hilltop hexes 0815, 0915, and 0916) and those hexes occupied by other artillery pieces or units.

\caseheading{7.6} If the most direct path between an artillery piece and a target unit can be traced through two hexes, one of which is a blocking hex and the other of which is not, the path is traced through the non-blocking hex.

\caseheading{7.7} Artillery pieces are friendly to the player controlling the last unit (not leader) to enter or pass through their hex.

Whenever the hex occupied by an artillery piece is occupied by an enemy unit, the artillery piece is captured and is flipped to its captured (back) side to indicate this. It is now friendly to the player whose unit captured it. Whenever the hex containing a captured artillery piece is occupied by a unit friendly to the original owner, the artillery piece is recaptured and is flipped to its front side. An artillery piece cannot be disrupted or eliminated, though it can be captured and recaptured any number of times during a game. The active player can use any or all friendly artillery pieces to bombard enemy units during his artillery phase.

\caseheading{7.8} Artillery pieces have no combat strength or Zone of Control (9.1), and they ignore both friendly and enemy ZOCs for all purposes.
